\documentclass[8pt]{extarticle}

\usepackage[utf8]{inputenc}
\usepackage[czech]{babel} % Správné české dělení slov a znaky
\usepackage[T1]{fontenc}

% --- GEOMETRIE: A4 papír s minimálními okraji ---
% Okraje nastaveny na 5mm, což je hranice, kterou zvládne 99 % běžných tiskáren bez oříznutí textu.
\usepackage[a4paper, margin=10mm, noheadfoot, landscape]{geometry} 

% Vynulování skrytých systémových mezer hlavičky, patičky a horního okraje
\setlength{\headheight}{0pt}
\setlength{\headsep}{0pt}
\setlength{\topskip}{0pt}
\setlength{\footskip}{0pt}

\usepackage{multicol} 
\usepackage{amsmath, amssymb, amsthm}
\usepackage{mathptmx} % Písmo Times - velmi kompaktní, klasické novinové sázející natěsno
\usepackage[protrusion=true,expansion=true,tracking=true]{microtype}

% Extrémní smrsknutí mezer uvnitř matematických vzorců
\thickmuskip=1mu \medmuskip=1mu \thinmuskip=1mu 

% ZÁBRANA ZALOMENÍ VZORCŮ (Matematické věci zůstanou striktně na jednom řádku)
\binoppenalty=10000
\relpenalty=10000

% Vynulování vnitřních mezer mezi odstavci a řádky
\pagestyle{empty} 
\setlength{\parindent}{0pt} 
\setlength{\parskip}{0pt}   
\setlength{\parsep}{0pt}

% Drastické zúžení řádkování a mezer mezi sloupci
\renewcommand{\baselinestretch}{0.72} 
\setlength{\columnsep}{1mm} 

\begin{document}
% Extrémně malé písmo (4.8pt) s adekvátním řádkováním (5.2pt) pro maximální hustotu
\fontsize{4.8pt}{5.2pt}\selectfont 

% Využíváme celou plochu A4 na šířku, takže 5 nebo i 6 sloupců je teď ideální
\begin{multicols*}{5} 


% =======================================================================
STATICKE SLOV: Holy Grail $Build \in O(n)$, $Query \in O(1)$, FKS: $c$-universalni system hes funkci, if kolize v s $Pr\leq c/m$, chceme $h$ najita a vyhodnocovat v $O(1)$, Lemma: pro $h \in_R H$ je $\mathbb{E}[\#\text{kolizi}] = (c\cdot n^2)/(2m)$,A: pro $m=\lceil n^2 \cdot c\rceil$ je $\mathbb{E}[\#\text{kolizi}] < \frac{1}{2}\rightarrow P[\text{h koliduje}]\leq P[\#\text{kolizi} \geq 2\cdot\mathbb{E}[\#\text{kolizi}]]\leq\frac{1}{2}$ dle Markovovy: Pro $X\geq0: P[X\geq k\cdot\mathbb{E}[X]]\leq\frac{1}{k}$, tedy $P[\text{bez kolize}] \geq \frac{1}{2}$ + Lemma o dzbanu $\rightarrow \mathbb{E}[\# \text{pokusu na volbu } h] \leq 2$, jedno overeni $h$ trva $O(n)$ pomoci prihradkoveho trideni.k B: pro $m=n$ je $\mathbb{E}[\# koliz] < \frac{c\cdot n}{2}\rightarrow$ $h$ najdeme v $c\cdot n$ krocich. Velikost vsech prihradek $Q_i$ je $O(n)$, $\Sigma_{i=[n]}\leq n + c\cdot\Sigma_i p_i^2 \leq n + c\cdot\Sigma(p_i + 2\binom{p_i}{2})$, protoze $p_i^2 = p_i + p_i^2 - p_i \leq p_i + 2(p_i^2 - p_i) = p_i + 2\binom{p_i}{2}$, pokracujeme $= n + $ $c\cdot\Sigma p_i + c\cdot\Sigma2\binom{p_i}{2}$, prvni sum je soucet prvku $\leq n$ a druha je $\#$ kolizi $f$, coz je $\leq \frac{cn}{2} \rightarrow$ cele $O(n)$. Pamet params $f \in O(1)$, $g_1...g_n \in O(n)$, tabulka indexovana $f \in O(n)$, tabulky $Q_i \in O(n)$, Cas $O(n)$ prumerne na volbu $f$ a vsech $g_i$, $Lookup \in O(1) w.c.$ DETER SLOVNIK: $O(n^k)\rightarrow O(n^2)$ trie hloubky $k$, cil $O(n^2)\rightarrow O(n)$;$(f,g)$ je $q$-dobra, pokud $x \rightarrowtail (f(x),g(x))$ je proste a $f$ ma na $S$ max $q$ kolizi. Lemma $(f,g)$ je $q$-dobra a $r\geq Ln+1$, pak $\exists a_0-a_{2^r-1}, (f',g'), q'$-dobry a $q' = \text{min}(n, \lfloor 2^{3-r} \cdot q \rfloor \cdot n)$. Algo: zvolim $r > Ln +3 \rightarrow 2^{3-r}<1/n$, $S\subsetneq\{{0,1}\}^w, w\leq2Ln+O(1)$, 1. krok rozrezanim $x$ mam $(f,g)$, ktera je prosta, kolizi max $\binom{n}{2} \leq n^2$, pote 2x lemma. Na konci musime pamatovat $a_i$ a $b_i$.Dk Lemma:Setridime radky, podle \# prvku $S_i$, Aby $x \in S_{<i}$ \& $y \in S_i$ kolidovali, tak $g(x) \oplus a_x = g(y) \oplus a_i$, $a_i \in_R [2^r] \rightarrow \mathbb{E}[\#NK] = |S_{<i}|.|S_i|.2^{-r}$. Markov:$\mathbb{E}=2$ pokusech mame \#NK max 2x vetsi. $\Sigma$pres$\forall$oznacime $*$. A: $2^{1-r}\leq 1/n, S_{<i}\leq n \Rightarrow *\leq \Sigma |S_i| = n$. B: Odstranime trivialni pripady $|S_i|=0 or 1$, scitame prefix $* = \Sigma_{i<t}\lfloor 2^{1-r}|S_{<i}||S_i|\rfloor$, vyjadrime$|S{<i}||S_i|\leq \Sigma|S_j|^2$, $d^2 \leq 4\binom{d}{2}$, $\Sigma|S_j|^2 \leq 4q$. Derandomizace: $\mathbb{E}[T]=P(\alpha)\mathbb{E}[T|\alpha] + (1-P(\alpha))\mathbb{E}[T|\neg\alpha]$. Binarni strom nad sloupci, vrchol je \# prvku ve sloupci, nebo intervalu sloupcu. VAN.EB STROM:$U=[2^w],w=2^k,hi(x),lo(x)$. $vEB(S,U)$:$\min$ alone,$\max$, prihradky pro $0..\sqrt{U}$ $P_i:=\{lo(x)|x\in S', hi(x)=i\}$, nad ni $vEB(P_i, \sqrt{U})$ a Sum str $vEB({i|P_i\neq\emptyset}, \sqrt{U})$.$Find, Succ$ zanorime do prihrad $P_i$, if $\neq\emptyset$, jinak v Sum.$Succ(i)$, z nej v $O(1)$ min. $Insert$ porovnam s min,max, pak $P_i.Insert(j)$, if$P_i = \emptyset$, tak $P_i.Insert(j) v O(1)$ a $Sum.Insert(i)$. $\forall$ operace v $O(LLU)$. Pamet a jeji inic. az $O(U) \rightarrow$ simulace na RAMu predem vynul pameti, metoda MIX, $M$ pole velke $|U|$, $I$ indexy inic. prvku a $N$ pocet, $X$ pole ukazatelu na $I$.$X[i]\rightarrowtail I$. INTERVALOVY VEB: listy char. vektor mnoziny $S$. Hloubka$=w$, $U=2^w$. Vnitrni vrch. 1/0 podle toho, zda podinterval je neprazdny. $Find$, $Pred$, $Succ$ v $O(LLU)$, pokud umime bin. vyhled. na list-koren. X-FAST:  kukacci hesing, pamet $O(n\cdot LU)$, vyhledavame binarne pomoci prefixu do hesovaci tab, listy spojene do $\leftrightarrow$ spoj seznamu, levy prazdny odkaz syna x ukazuje nejlevejsi list v pravem podstromu x $\rightarrow$ queris $O(LLU)$, $Update$ v $\overline{O}(LU)$. Y-FAST: Indirekce $\rightarrow$ lepsi pamet a updates, $n$ prvku na bloky $B=\Theta(LU)$, uvnitr bloku bvs, nejvetsi prvek v bloku je reprezentant, tech je $O(n/LU)$, Pamet za $\forall$bvs $O(n)$ + pamet za x-trii $O((n/LU)\cdot LU)=O(n)$. $Update \rightarrow Find + Ins/Del_{bvs} \in O(LLU)$, kazdy blok $\rho\in[1/2B, 2B]$, pokud pre/podtece, tak delime/slucujeme a ins/del v x-trii, $O(LU)$, ale nejmene $\Omega(LU)$ $Ins/Del$, proto $\overline{O}(1)$. $Del$ reprezentanta $\rightarrow$ pomnicek! HALDA:$Ins,DecreaseKey\in O(L_{k}n)$, $ExtractMin\in O(k\cdot L_{k}n)$.Dijkstra $n\times Ins + m\times DKey +n \times EMin$ nabyva min pro $k=m/n$. BINOMIALNI H:Les Binom. stromu. $B_0 :=$ 1 vrchol, $B_k$ ma deti $B_{0..k-1}$, $k+1$ hladin, $2^k$ vrcholu, max $k$ deti. Strom radu $k$ max 1$\rightarrow$ jednoznacne urceno ktere. $\#$strmu $O(Ln)$ Repre v pameti: vrchol$:=$(prvek, priorita, ptr na rodice, ptr na nejlevejsi dite, ptr na left a right sibling). $Merge$ jak scitani dvou bin cisel$\in O(Ln)$, $Ins$ je $Merge$ s $B_0$. $Build$ v $O(n)$ misto $O(n\cdot Ln)$. $EMin\rightarrow$ $Find$ + $Merge \in O(Ln)$. $DKey \in O(Ln)$. $Del\rightarrow DKey(-\infty) + EMin$. LINE BIN H: Porusime invariant max 1 strom od radu $r$, $Merge \in O(1)$ pomoci obousmer spojaku korenu,$Insert \in O(1)$ prilepime do spojaku, $Build$ pro $n$ prvku $\in O(n)$, behem $ExtractMin$ konsodilace ta v $O(Ln + \#\text{stromu})$, $\Phi := t =\#$ stromu, pri $EMin$ klesne na $O(Ln)$, Amort Cena = Real Cena + $\Phi_\Delta$, $\Phi_\Delta \leq Ln - t =$ po minus pred, Real cena $EMin = O(Ln + t)\rightarrow$ Amort cena $=(Ln + t) + (Ln - t) \in O(Ln)$, konsodilace pomoci prihradkoveho trideni s $Ln$ prihradky. $Decrease \in \overline{O}(Ln)$ i $\in O(Ln)$, nemeni $\Phi$. $Delete$ a $Increase \in O(Ln + t)$, ale $\in \overline{O}(Ln)$. FIBONACC HALD: Bily vrchol: Neztratil dite nebo je koren, Cerny: ztratil jiz jedno dite. Vrchol v pameti ma: prvek, jeho prioritu, rad (\# deti ), flag bily/cerny, ptr na rodice, ptr na prvni dite, ptr na left i right sourozence + si udrzuji global min cele haldy. Fib cisla $F_0=0, F_1=1, F_{n+2}=F_{n+1}+F_n$ rostou $\Theta(\varphi^n), \varphi=1.618$, Lemma0 $F_0+F_1\ldots F_k=F_{k+2}-1$ (indukci). Lemma* vrchol radu $k$ (ma $k$ deti) ma aspon $F_{k+2}$ potomku, vcetne ($T(u):=$ velikost podstromu $u$). Dukaz: $\odot$ v konsodilaci ze 2 stromu radu $k$ vznikne str radu $k+1$. Indukci dle hloubky, vrchol $u$ s detmi $x_1,..,x_k$ a jejich rady $r_1,.,r_k$, $\odot$ poradi vrcholu v jakem byly stromy slucovany. Tvrzeni $\forall i: r_i \geq i-2$. $T(u) \geq 1 + \Sigma_{i=1\_k}(T(x_i))$ (deti maji rady $r_i$, tedy dle IP) $\geq 1 + \Sigma_{i=1\_k}(F_{r_i+2})$ (pouzijeme Tvrzeni a $T(x_1) \geq 1 = F_2$) $\geq 1 + F_1 + F_2 + .. + F_k$ (pripiseme plus $ F_0 = 0$ a dle Lemma0) $=1 + (F_{k+2} -1) = F_{k+2}$. Dk Tvrzeni: Pro $r_1, r_2$ plati. Okamzik, kdy $x_i$ pripojeno pod $u$. $u$ musel mit aspon $i-1$ deti, mohl jich mit vic, protoze nektere jsme mohli pozdeji smazat, tedy rad $u \geq i-1$. Spojujeme stromy se stejnym radem, proto $r_i \geq i-1$, protoze $x_i$ neni koren, mohlo prijit od te doby max o 1 dite a stat se cernym $\rightarrow$ rad $\geq i-1-1 = i-2$. Diky tomu vsemu je max \# deti (tedy rad) vrcholu v $O(Ln)$. $\Phi = \#$ stromu $+ 2\cdot\#$ cernych vrcholu. $Cut(x)$: barva($x$)$\leftarrow$bila, odstranime hranu $rodic\rightarrow x$, if $u$ koren tak done!, jinak podle barvy ho obarvime na cerny nebo uz byl a $Cut(u)$. $Decrease(x)$ - mrk na rodice a if porusena priorita $Cut(x)$, wc slozitost $O(1 + \#cernych)$, amortizovane: kdyz kaskadujeme, tak menime cerny $\rightarrow$ bily $=-2$ a plus strom $=+1$ $\Rightarrow \Phi_\Delta = -1$ a zaplatime tim kaskadovani, a zastavime na bilem, ze ktereho udelame cerny a tim se zvedne potencial o konstantu $\rightarrow \overline{O}(1)$. Operace $Ins, Merg, Build$ same jak u linych (i amort). $EMin$ - wc $\in O(Ln + \#\text{stromu})$, $\overline{O}(Ln)$, stejne jak u linych hald. $\Rightarrow$ Dijkstra bezi v $\overline{O}(n\cdot 1 + n\cdot Ln +m\cdot1)$. SPLAY STROM: $T(v)=\#\text{vrcholu}$, $s(v)=|T(v)|$, $r(v)=Ls(v)$, $\Phi(v)=\Sigma_{u\in T(v)} r(u)$. $sr\Phi$ nabyvaji hodnot $1\_n$, $0\_Ln$, $0\_nLn$. Thm Access: Amortizovana cena $Splay(x) \leq 3\cdot(r'(x)-r(x)) +1$ ($\in \overline(O)(Ln)$) ($r'$ je po operaci $Splay$).Operace $O_i$ ma ceny $R_i$, $A_i$. Dale $A_i = R_i + (\Phi_i -\Phi_{i-1})$, kde $\Phi_i$ je potenc po operaci $O_i$. Uteleskopujeme po $m$ operacich na $\Sigma R_i = \Sigma(A_i) + \Phi_0 - \Phi_m$. Vyhodime $\Phi_m$, ale $\Phi_n$ je az $O(nLn)$, proto $m \times Splay$ trva $O(mLn + nLn)$, (chceme \# operaci aspon \# vrcholu). Lepsi analyza: vahy $w(v)>0$, premapujeme $s(v)=\Sigma w(v)$, proto $r$ a $\Phi \in (-\infty, +\infty)$. $\odot$ThmAccess plati, protoze nerovnost pro logaritmy (konkavnost), po odkrajeni podstromu velikost, ranky, $\Phi$ nestoupnou. Neplati $(r'-r) \in O(Ln)$. Parametrizujeme: $3(r'-r)+1 \in \overline{O}(r_{max} - r_{min} + 1)$. $Find$ stejne jako normalne $\rightarrow\in\overline{O}(Splay)$. $Insert$ nemuzeme pridat do listu a pote $Splay$, skrz pridani vrch s velkou $w$, skokove vzroste $\Phi$. Prvni $Splay$ naslednika $x$ a pote podrozdelit levou hranu (levy podstrom $<x$, pravy $>x$). Zmena $r$ $O(1)$ vrcholum, cas $\overline{O}(Splay)\cdot O(1)$. $Delete$- $Splay$ vrcholu, pak ho smazeme a znovu $Splay$ extremu a prepojeni, potencial muze jen klesnout, takze stejne jak $Insert$. LemmaW: Amort cena $Splay(x) \in O(L(W/w(x))+1)$. $r_{max} = LW$. Dk: $Splay \leq O(r_{max}-r{min}+1) =$ $O(LW - Lw(x) +1) = O(L(W/w(x))+1)$. Rovnomerne rozdeleni (tedy puvodni Splay Strom, neuprednostnujeme zadny prvek): $W=n\cdot\frac{1}{n}=1$. $\frac{1}{n}\leq s\leq1$,$-Ln\leq r\leq0$, $-nLn\leq\Phi\leq0$. Dosazenim do LemmaW $Splay\rightarrow O(Ln)$. Celkova cena $m$ operaci: $O(mLn)$ za amortizovane operace, musime pricist $\Phi_0-\Phi_m \leq 0-\Phi_m\leq-(-nLn)=nLn$, tedy $m$ operaci stoji $O(mLn + nLn)$. $\odot$ vahy muzeme vzdy normalizovat, aby $W=1$, vynasobime totiz $\forall s$ konstantou $c$, tedy $r = L(c\cdot s(v)) =Lc + Ls(v)$, pouze se k $\forall r$ pricte $Lc$, tedy $r'-r$ je v poho.| Staticky opt str: klice $X={x_1\_x_k}$, frekvence $f(x)$ nebo $p_i$, prst ze $x_i$ bude dotazovan. bvs strom $T$, postavime v $O(n^2)$, $c_T(x)$ cena pristupu k $x$ = \# vrcholu na ceste $x$-koren. Thm(StaticOpt): $\varphi=x_1\_x_m$ pslpnst pristup na $X$: $\forall x$ se vyskytuje $f(x)$-krat, $f(x)>0$. $T$ je stat opt str, pak cena $Splay$ na $\varphi$ je $\in O(c_T(\varphi))$. Dk: $w(x) = 3^{-c_T(x)}$; ukazem $W<1$, pro nekonecny bvs $\Sigma_{i\geq1} 2^{i-1}\cdot3^{-c_T(x)}=\Sigma(2^{i-1}/3^i) = (1/2)\cdot\Sigma(2/3)^i = 0.5(-1+\Sigma_{i\geq0}(2/3)^i)$, nekonec rada $\Sigma = 1/(1-(2/3))$, cele $=1$, pro konecny $<1$.LemmaW $Splay \in O(L(1/3^{-c_T(x)}))$ $\in O(c_T(x)+1)$, $c_T\geq1\rightarrow O(c_T(x))$; Musime pripocitat $\Phi_0-\Phi_m$: $3^{-c_T(x)}\leq s(x)\leq 1$, $-c_T(x)\cdot L3\leq r(x)\leq0$, $\Omega(\Sigma c_T(x)) \leq \Phi \leq 0$, proto $\Phi_0 \leq 0$ zahodime, $-\Phi_m \in O(\Sigma_{x} c_T(x))$, protoze na kazdy prvek se ptame aspon 1, tak $\in O(c_T(\varphi)) \square$. Dynamicka verze otevreny problem.| Working Set: hledame $x_i$, ws $z_i$ je mn prvku, ktere jsme hledali mezi tim, pokud jsme $x_i$ jeste nehledali, tak $z_i$ jsou $\forall$ prvky od zacatku hledani. Cil ukazat amort cena $Splay$ je $O(L(1+z_i))$. Thm(WorkSetBound): Cena pristupu $x_1\_x_m$ je $O(nLn+m+\Sigma_i L(1+z_i))$. Dk: lru list, prvky jsou usporadane zleva, posle last pristupu, $\odot$ hledame $x_i$, pak jeho pozice $p(x_i)$ v lru je jeho $z_i$. Nastavime $w(x)=1/(p(x)+1)^2$. $\odot \Sigma(1/k^2)$ konverguje $\rightarrow W \in O(1)$, dosadime do LemmaW $Splay$ stoji $O(L(p(x)+1)^2+1) \rightarrow$ pres $m$ operaci $O(m)$ (za +1) plus $O(\Sigma_i L(p(x_i)+1)) = O(\Sigma_i L(z_i+1))$;Zmena vah: po $Splay(x)$ je $p(x)=0$ tedy $w(x)=1$ a vzroste z puvodni nejhur $w(x)=1/m^2 \rightarrow \Phi_\Delta \leq 1$, ostatnim $y$ vzroste $p(y)$ o 1, takze snizi $\Phi$, proto nam to nevadi; opet musime pripocitat pro $m$ operaci $\Phi_0 -\Phi_m$: $1/n^2 \leq w \leq 1$, $1/n^2 \leq s\leq W$, $-2Ln\leq r\leq O(1)$, $-2nLn \leq \Phi \leq O(n)$, $\Phi_0 - \Phi_m \in O(nLn) \square$ STROMOVE STRUKTY: Drzi les s ohodnocenim V/E, Operace: spojeni stromu, $\exists$ cesta mezi dvema vrcholy, min/max na ceste, updaty cestove ($+\delta$ vsem hranam), bodove ($+\delta$ hrane). DS CESTA, staticka, $n$ ($=2^k$) vrcholu, ktere maji ohodnoceni ($w$), postavime intervalovy strom, vnitr vrcho pamatuji min, interval dotaz lze pokryt kanonyckymi podstromy tech je $O(\text{hloubka})=O(Ln)$, updt bodovy - staci zmenit list a pak minima na ceste korenlist, updt cest - pridam znacku $+\delta$ do vsech korenu kanonickych podstromu intervalu; kanonic podstromy najdu pomoci LCA(x,y) a pote to jsou $\forall$ vrcholy uvnitr.  DS STATIC STROM (HEAVY-LIGHT): hrany orientovane ke koreni, hrana $u\rightarrow v$ tezka $\equiv |T(v)|>\frac{1}{2}|T(u)|$, $\odot \forall$ vrchol $u$ ma max 1 tezkou hran z deti $\rightarrow$ tezke disjunkt cesty, projdu-li lehkou hranou nahoru $\#$ vrcholu v podstrom se aspon zdvojnasobi $\rightarrow$ \# lehkych  i tezkych projdu po ceste do korene max $O(Ln)$. $Build$ v $O(n)$ - 2$\times$dfs, 1. velikosti podstromu, 2. poradi vrcholu na tezk ceste, rekurzi na lehke. Vrchol - kolikaty na tezke c a id tezke c$\rightarrow$ $lca$ jdu po lehkych hranach, skacu po tezkych cestach, dokud nenajdu stejnou $O(Ln)$. Tezka c $=$ intervalovy str, $cquery(xy)$ - lca(xy), pro $\forall$ navstivene tezke c interval dotaz (suffixy tezkych cest, krom last), $O(Ln)$ tezkych c a lehkych hrn $O(Ln + L^2n)=O(L^2n)$, $hupdt = O(Ln)$ (pouze pro jednu tezkou cestu $hupdt$) a $cupdt = O(L^2n)$ jako $cquery$. LINK-CUT STRMY (DYNAM VERZE) - hrany tluste/tenke, opet max 1 tlusta $\rightarrow$ tluste c, interface $Parent(v)$, $Root(v)$, $Cut(v)$, $Link(u,v)=$ $u$ pod $v$, $u$ musi byt koren, $Evert(v)=$ prekoreni na v, $Cost(v)$, $CestaMin(v)=$ min z v-koren $SetCost(v,\alpha)$, $CestaUpdt(v,\delta)$, $Expose(v)$.  Tlst cesta $=$ Splay strm, in-order = poradi na ceste, $\forall v$ pamatuji $min$ z podstrm, $+\delta$. Interface tlstych cest: $Prev, Next, First, Last$ - (first = stale doleva) +splay; $Cut =$ splay + del hrany vpravo pod $v$; $Link$ = $Ssplay(Last(v))$ a pripojim vpravo; $Cost, SetCost$; $CMin =$ od $v$ do konce cest$\rightarrow$ splay(v) a $min$ z praveho podstromu; $Cupdt =$ Splay(v) + liny updt $+\delta$ do $v$,$\odot$ po Splay $v$ nema v sobe jine $+\delta$, vycistime pred rotacemi, jdeme z $v$ do korene pak zpet; $Reverse =$ opet line, pri pruchodu ze shora dolu cistim, jinak pri splay nevadi absolutni orientace, staci relativne (zigzig, zigzag) a cistime jen dve nad. Operace nad Tlst cest $O(Ln)$. Odkaz na tenkeho rodice tlst cesty ma koren tlst cesty, pri splay predelame do aktual korene.|Spolecny strm: $Expose =$ pomoci $Link$ a $Cut$ propojim tluste cesty z $v$ do korene a udelam jednu tlstou cestu. 1. musim vySplayovat vrcholy pod kterymi jsou tlst cesty do korenu ve splay stromech,2. pote $Link/Cut$ a muzu dobre pripojit podstromy abychom meli jeden velky splay strm,3. nakonec v nem $Splay(v)$. Line updaty v $Splay$ strmech jsou v poradku skrz $Splaye$, $min$ podstromu ve vrcholech u prvnich $Splayu$ ok, pri $Link/Cut$ musime projit nove vzniklou $abc$ cestu a prepocitat od spoda nahoru, posledni $Splay(v)$ ok. $\Rightarrow$ $\forall$ global operace pomoci $Expose \cdot O(1)$, napr $Evert(v)\rightarrow Expose + Reverse$ ve vyslednem splay strome.| Amort $Expose$: $s(v) = \#$ potomku ve spolecnem stromu, $1\leq s\leq n$, $0\leq r\leq Ln$, $0\leq \Phi\leq nLn$, $\odot$ prohozeni lehke a tlste hrany ho nemeni $\rightarrow O(1)$, 1. najit $abc$, pak $Splay(a), Splay(b).. \rightarrow O(r'-r+1)$ (dle AccessThm) $\rightarrow O(r'(a)+r(a)+1)+O(r'(b)+r(b)+1)$ (uteleskopuje skrz $r(a) \geq r'(b)$ a \# vrcholu $abc \times (+1)$) $\rightarrow O(r'(a)-r(v)+\# abc)$, 2. pri $Link/Cut$ $O(1)$ za kazdy $abc..$ vrchol, 3. $Splay(v)$ zaplati $O(\# abc)$ z 1. a 2. $\Rightarrow Expose$ v $\overline{O}(Ln)$, wc je $O(n)$. Aplikace LinkCut: Dinic $O(n^2\cdot m)$, rezervy jsou vahy, vrstevnata sit zorientovana doprava je LinkCut strom. DYNAMIZACE: 1. Plna prestavba - kdy? (kapacita, parametry) 2. Castecna prest.: Pro bin strom: $v$ ma syny $l$ a $p$, plati $|T(l)| \leq \alpha|T(v)|$ pro $\frac{1}{2}\leq\alpha\leq1$ $\rightarrow$ hloubka $O(L_{(1/\alpha)}n)$, udrzujeme pomoci pocitadel, pri $Insert$ pridam +1 na ceste do korene, kontroluji podm, kde nejvyse porusena $\rightarrow$ prestavim podstrom (in-order pomoci dfs, do pole a medianuji), $\Phi = ||T(l)| - |T(p)|| \rightarrow O(Ln)$ na $Ins$ i $Del$. DYN 2D INTERVAL STR: bb-$\alpha$ vyvazovane strmy: Primarni strom dle $x$, listy obsahuji prvky, vrchol ma - delici klic $x_{mid}$, interval jeho podstromu $x_{min}$, $x_{max}$, ptr na sekundarni strm serazeny dle $y$ a obsahujici $\forall$ prvky, ktere $p_x \in[x_{min}, x_{max}]$; $Build$ - setridime prvky dle $x$ + dle $y$ ($O(nLn)$), rekurzivne: pro vrchol $v$ vezmeme median pro jeho interval z setridenych $x$, vyrobime sekundarni strom $\rightarrow$ rekurze na syny, $Build$ sekundarniho v $O(\#\text{prvku co do nej padly}) \rightarrow \Sigma$ pres hladinu $O(n) \rightarrow$ celkove pres $O(Ln)$ hladin $\Rightarrow O(nLn + nLn)$. $RectangleQuery$: return \# prvku v obdelniku, v $O(L^2n)$ skrz kanonicke podstr. Prostor $O(nLn)$ - 1 vrchol lezi v $y$ strmech na ceste z neho do korene v $x$ strme. Dynamizace line: pri rekurzy zpatky po $Ins$ kontrolujeme, zda $T(syn) >\alpha T(v)$, if ano, prestavime podstrm velikosti $k$ v case $O(kLk)$, $\#$ operaci od posledniho $Rebuild$ je $\Omega(k) \rightarrow (\text{Cena} Rebuild)/\Omega(k) = O(Lk)$, prvek $v$ musi predplatit $Rebuild$ pro $\forall$ hladiny, proto $Ins$ nebo $Del$ $\in O(L^2n)$. OBECNA DYNAM: Vyhledavaci Problem $f: U_Q \times 2^{U_X} \rightarrow U_R$. Vyhled Prb je $rozlozitelny \equiv$ $\exists \sqcup: U_R \times U_R \rightarrow U_R$ vycislit v $O(1)$, t.z. $\forall A,B \subset U_X, A\cap B = \emptyset, \forall q \in U_Q$ $f(q, A\cup B) = f(q, A) \sqcup f(q, B)$. 1. Zda $v\in S$? ano, $\sqcup$ je bool or, 2. Nejblizsi bod $s\in S$ bodu $v$? ano, $\sqcup$ je min z dvou min, 3. $v\in conv(S)$? ne. Mame statickou ds: $B_S(n)=$ cas buildu, $Q_S(n)=$ cas dotazu, $S_S(n)=$ pamet, navic $B/n$, $S/n$ a $Q$ jsou neklesajici v $n\rightarrow +\infty$. Bloky $B_i \equiv$ dvojkovy zapis $n$ $\rightarrow O(Ln)$ bloku. Semi-Dynam ds: $Q_D(n) = O(Ln\cdot Q_S(n))$ (diky neklesavosti), $S_D(n) = \Sigma_{i\in I}S_S(2^i) = \Sigma(S_S(2^i)\cdot2^i/2^i)$ $\leq S_S(n)/n \cdot \Sigma2^i \leq S_S(n)$. $Ins$ - najdeme nejblizsi prazdny blok $B_0$, tedy $B_i$, rozebereme vsechny po $i-1$, pridame k nim prvek $x$ a vybudujeme $B_i$, cas:$B_S(2^i)$, muselo probehnout $2^i$ operaci $\rightarrow \leq B_S(2^i)/2^i \leq B_S(n)/n =$ prace na 1 hladinu pro 1 prvek, $Ln$ hladin a $\forall$ prvky z jednoho radu do dalsiho max 1x $\rightarrow \overline{I_D(n)}=O(B_S(n)/n \cdot Ln)$. Priklad sorted array: $B_S\in O(nLn)$, $Q_S\in O(Ln)$, $S_S \in O(n) \rightarrow$ $\overline{I_D(n)}\in O(L^2n)$.. WC $I_D \in B_S$: 3 hotove $B_i$ a 1 rozestaveny $B_i$,$\odot S_D(n) = O(S_S(n))$, wc $Insert \in O(Ln + \Sigma (B_S(2^i)/2^i))$ (rezie za hladinu + pres hladinu $1/2^i$ prace) $= O(Ln + B_S(n)/n \cdot Ln) = O(B_S(n)/n \cdot Ln)$. Plna dynam pomoci pomnicku, mame stejny (amort) rozklad na $B_i$, $Delete$ - z prvku ptr na blok a tomu reknu a bloku reknu "udelej pomnicek"; adresar, ktery si pamatuje pro prvek $\rightarrow$ blok, bude to bvs, dotazy a operace $O(Ln)$, pro radove $\Theta(n)$ operacich prebudujeme. PERZISTENCE efemerni, semi-perzist, plne perzist, funkcionalni (objekty se nikdy nemeni). Funkc seznam s prependem, linked list, pridavame na zacatek, verze ptr na vrchol, $O(1)$ na $Ins$, $Del$ i prostor na verzi. Funkc bvs: pridame vrchol $x$, musime zkopirovat rodice a zmenit jeho ptr na $x$, opakujeme az do korene; verze ptr na novy koren, $O(Ln)$ na verzi, $Ins$ i $Del$ to nezpomali. 
$\odot$ pri vyvazovani se $\forall$ zmeny deji na ceste do korene $\rightarrow$ avl stromy s kopirovanim cest, zpomaleni $O(1)$-krat, prostor/verze $O(Ln)$. Aplikace Lokalizace bodu v Polygon rovine: 1. reseni: pasy horizont, dle udalosti $rightarrow$ jeden pas rozdeluji usecky, ty seradime, $Build$ a prostor $\in O(n^2)$, cas dotazu $O(Ln)$. 2. poridime bvs, ve kterem budeme pridavat/ubirat usecky v bodech udalosti od shora dolu, vyhledavani: vyhledame verzi binarne a pote v bvs, cele $\in O(Ln)$; zlepsi prostor a $Build$ na $O(nLn)$. Nevyhoda funkc bvs, nemaji ukazatele na rodice. POINTER DS: krabicka=($O(1)$ slov dat, $O(1)$ ptrs), Drzatka $O(1)$ ptrs na krabicky. Bvs je priklad ptrve ds. Chceme zmenit krabickovou ds na semi-perzist. Musime verzovat: -data a ptrs v krabickach (slouzi k $Ins$ i $Query$ - strukturalni zmeny), -nemusime flags a pomocne data, ktere , napr u avl znamenka/hloubky (nestrukturalni - slouzi jen $Ins$, ne $Query$). Cil strukt. zmena $\in \overline{O}(1)$ a prostor taky $\rightarrow$ (2,4)-strom $\rightarrow$ amortizovane $O(1)$ strukturalnich zmen na operaci, cas $O(Ln)$ na operaci, $\rightarrow$ tedy nase masinerie na semi-perzist zvladne amortizovane cas $O(Ln)$ na operaci toho (2,4)-stromu.| Technika Tluste Vrcholy: krabicka: (pocatecni vrchol,(.. (verze, novy obsah))), $O(1)$ prostoru a casu na vytvoreni verze/zmenu, dotaz $\rightarrow$ binarne najdeme verzi $O(Lh)$, $h=$delka historie, pouzijeme $y-fast trii$ pak $\overline{O}(LLh)$. Problem omezeny \# verzi. Obecna semi-perzis DS: Prepoklad $deg^{in}(v) \leq p$, $p=$konstanta. Krabicka: (pocatec ver, data, ptrs, $e$ slotu na zmeny, zpetne ptrs), $e\geq p$; pridavam zmenu "v teto verzi $i$ menim ptr $j$ (poradove cislo) na ptr $P$(samotny ptr)", if v krabce misto, hotovo. Pri $Query$ prochazim ver v krab, aplikuju a skoncim na hledane ver $\in O(1)$. Pri preteceni cp krab. $\odot$ zpetne ptrs nemusime ver. Cp data, ptrs, set new ver + prenastavit krabky ukazujic na starou skrz zpet ptrs $\rightarrow$ prida zmeny do jejich zurnalu $\rightarrow$ dalsi cp. Aby se cp zastavilo "$\forall$ ptrs zmenime v jedne ver max 1x", pokud menime ptr u rodice, ktery ma "pocatec ver"=aktual, tak menime ptr primo a neverz zmenu, nebo se podivam, jestli v nejakem ze slotu na zmeny neni zaznamenana zmena toho ptru v aktualni verzi a poku ano, tak prepisu primo tento zaznam zmeny $\rightarrow \forall$ krabky cp max 1x v jedne verz. Amort: $\Phi= \Sigma$ zaplnenych slotu (zmen) z krabek dosazitelnych z aktualni ver, $\varphi(\text{krabky} v)= \#$ plnych slotu $\leq e+1$ (if plna krabka a pridame zmenu, jeste zapocitame). Copy krabky $v$: potencial $e+1 \rightarrow 0$, $+1$ vyuzijeme na kopirovani krabicky, uvolnene $e$ vyuzijeme na projiti $p$ zpetnych ptru a zapsani zmen do rodicu, vyjde, protoze $p\leq e$ $\Rightarrow$ cas i prostor strukturalni zmeny v $\overline{O}(1)$. Reprezentace krabky: cislo pocatecni verze, obsah zakladni verze: data+ptry, ptr na nejblizsi novejsi kopii (spojak z verzi), $p$ zpetnych ptru, $e$ slotu na zmeny, verze: (cislo verze + hodnoty vsech drzatek (muzeme je vzdy copy, protoze jich je $O(1)$)). Ptr na nejblizsi kopii, je proto ze kdyz menime v nejake verzi $z$ zpetne ptry, na aktual krabku, mobla byt krabka na kterou sahame skrz zpetne ptry zkopirovana, takze musime sahnout do te nejblizsi kopie, ale to max 1x v aktualni verzi.| Plna Perzistence: Musime vychazet z neceho w.c., protoze se muzeme pridavat vzdy do verze, kdy nastala draha operace a to pokazi amortizaci. Vsechny zmeny jsou strukturalni, (musime verzovat i znamenka/atd). Problm jak porovnavat verze $\rightarrow$ linearizace historie, prochazime strom, pri zanoreni aplikujeme zmenu pri vynoreni jeji inverzi. Navic vyzadujeme $deg^{out} \leq d$ ($d$ konstanta). Zbytek slozity. LIST LABELING: list order - pslpnst prvku, $Insert(z,x,y)=$ $z$ mezi $x,y$. ListLabeling $\rightarrow$ $\forall$ prvky maji label, ktery udava poradi. Exponenc ll: \# prvků $\leq M$. Zarážky jsou na pozicích $0, 2^M$. Nový prvek dám na průměr sousedů. Po $M$ insertech mezery aspoň 1. $O(1)$ na insert w.c. Funguje pouze pro $M \in O(w)$, kde $w$ je šířka slova. Polynomialni ll: bin. strom hloubky $O(\log M)$. Externí vrcholy jsou prvky. Kroky ve stromě L/P jsou bity v labelu (dopadinguju). Dřivější prvky dostanou nižší labely. \# bitů labelu je $O(\log M)$. $Insert$ znamená změnit příslušný externí vrchol na interní a pod něj připojit přelabelovaný prvek a nový prvek. Vyvažovat budeme líně s BB-$\alpha$ stromy. Přebudování stojí $O(k)$ (velikost podstromu), kde $k$ je \# prvků podstromu a musíme je přelabelovat, to se ale schová do $O(k)$ a tedy amortizovaně insert běží jako v líně vyvažovaném, tedy $O(Ln)$.| Kombinace ll: Rozdělíme posloupnost na bloky velikosti $\theta(\log n)$ Uvnitř bloků používáme labelování v exponenciálním rozashu. Nad bloky používáme labelování v poly rozsahu. Prvek má (label bloku, label uvnitř bloku). Dokud $Insert$ do bloku, tak v $O(1)$. Jakmile přeteče, tak ho rozdělíme na dva poloviční (v kazdem rozdrobime prvky rovnomerne) a přidáme blok v globálních blocích, tedy musíme přelabelovat všechno uvnitř jednoho bloku. Čas $O(\log n)$ na přelabelovaní, čas na globální $Insert$ $O(\log n)$ amortizovaně. Amortizovaně $O(\frac{1}{\log n})$ krát. Zrada, musí se přidat indirekce, kvůli přelabelování ostatních bloků při velkém globálním rebuildu, ve kterých je $O(\log n)$ prvků, jinak by to bylo $O(L^2n)$. CACHE STRUKTURY: Van-Emde-Boasovo uložení binárního stromu: Uložení úplného bin. stromu s hloubkou $h=2^k$: Rozřízneme v $h/2$. Uložíme nejdříve vrchní strom, poté těch hodně nižších zleva doprava. Pro paždý podstrom se rekurzíme. Zascopujeme na hloubku rekurze $i$, kde už se $S_i \leq B$, nad touto hloubkou rekurze je $S_{i-1} > B$ a taky $S_{i-1} = S_{i^2}$. Z toho plyne, že na této hloubce je $\sqrt{B} < S \le B$.
 Hloubka $S$ je tedy $\frac{1}{2}\log B < \log S \le \log B \implies \log S \in \Theta(\log B)$.
 Každý strom v tomto scopu se načte do cache v 1 až 2 operacích $\implies O(1)$ práce. $\#$ kusů je $O(\frac{\log n}{\log B})$. $\implies$ cesta z kořene do listu stojí $O(\frac{\log n}{\log B}) \cdot O(1) = O(\frac{\log n}{\log B})$. Uložení úplného bin. stromu s obecnou hloubkou $h$: Problém je, že při určování hloubek horního a spodních stromů musíme zaokrouhlovat nahoru nebo dolů (+-1) $\rightarrow$ může se posčítat blbě. Trik: $h_b :=$ nejbližší mocnina 2 vzhledem k $h$, $h_t := h - h_b$. Lze ukázat, že poté pouze první strom, úplně nahoře bude anomálně velký a jinak ostatní už pěkně mocnina 2. PACKED MEMORY ARRAY: Struktura: Mám pole s mezerami, rozdělím na kusy velikosti $\Theta(\log n)$, kdy $\#$ kusů = $2^h$, $h$ je hloubka bin. stromu, který nad kousky stavím. Hustota intervalu $:= \frac{\text{\# prvků}}{kapacita}$, standartní hustota v hloubce $i$ je $\in [\frac{1}{2} - \frac{i}{4h}, \frac{3}{4} + \frac{i}{4h}]$. $Insert(x)$:  Přidám ho do kousku, pokud přeteče, najdu na cestě kousek-kořen první vrchol, který nepřetekl přes svou std. hustotu. $\forall$ prvky přerozdělím rovnoměrně pod tímto vrcholem. Pozorování, některé vrcholy nad tímto vrcholem budou stále přeteklé nad std. hust., ale to nevadí, protože hustota vrcholu je aritmetický průměr potomků na jedné hladině. Pokud dojdu až do kořene, proběhně celostromová reorgranizace.| Amortizovaná cena $Insert(x)$: Provádíme reorganizaci vrcholu, který $\rho \leq \frac{3}{4} + \frac{i}{4h}$ a $\exists$dítě $d: \rho_d > \frac34 + \frac{i}{4h}$. Pokles hustoty aspoň o $\frac1{4h}$. Pokud má rodič kapacitu $k$, tak ubude aspoň $(\frac1{4h})\cdot(\frac{k}2) = \frac{k}{8h} \circeq  \frac{k}{\log n}$ prvků. 
 Přeorganizace intervalu délky $k$ trvá $O(k)$, proto amort. cena $Insert(x)$ je $O(log^2(n))$, protože ještě musíme sečíst přes všechny hladiny.| Reorganizace celého pole: Musíme zvětšit pole tak, aby hustota v novém kořeni vyšla mezi $[\frac12, \frac34] \rightarrow$ mezi další globální reorg. musí proběhnout lineárně mnoho $Insert(x)$. Počet nových kousků musí být mocnina 2.|
 Jak počítat hustoty? Spočítáme hustotu, tedy počet prvků, jednoho nalezeného kousku\dots $\Theta(\log n)$ Přidáme levý/pravý sousedící interval dvakrát delší délky. Opakujeme, dokud nepokrejeme interval velikosti $k$. Stojí to $\Theta(k)$, reorganizace těchto prvků stojí $O(k)$, proto se to zaplatí.| I/O složitost: počítání hustot a reorganizace: Čtení kousků, pro počítání hustoty jde pomocí dvou skenů, jeden $\leftarrowtail$ a druhý $\rightarrowtail$, kdy ty stále dvojnásobné intervaly, které přidávám vpravo čtu pomocí $\rightarrowtail$ a vice-versa.
Reorganizace: 1. zťukám prvky doprava v intervalu, který reorganizujeme, pomocí dvou skenů, kdy jeden čte, druhý zapisuje. 2. rozházím zťukané prvky pomocí dvou skenů, kdy oba stejný směr a jeden čte a druhý zapisuje. Nikdy nenarazím na zťukanou část. Max šest skenů $\Rightarrow \Theta(\frac{k}{B} + 1)$ Amortizovaná I/O složitost bude $O(\frac{\log^2 n}{B} +1)$, protože Insert stojí amortizovaně tolik a I/O operací je $B$ krát méně.| Udržování vEB bin. stromu nad PMA: Musíme přepočítat klíče v předcích prvků, které reorganizujeme. Celková cena $O(\frac{\log n}{\log B} + \frac{k}{B} + \frac{\log k}{\log B})$.
$\frac{\log n}{\log B}$ za cestu z $LCA$(krajů reorganizovaného intervalu) do kořene, to ale stejně potřebujeme na nalezení kusu, do kterého budeme vkládat. $\frac{k}{B}$ za všechny stromečky, co jsou uvnitř, vleze se do $O(\frac{k}B +1)$, tedy amortizovaně $\frac{\log^2 n}{B} + 1$. $\frac{\log k}{\log B}$ za stromečky, co jsou na krajích, za každou hladinu max dva, to se ale vleze do $O(\frac{k}{B} + 1)$. Celková amortizovaná složitost $Insert$ včetně udržování vEB bin. stromu nad je $O(\frac{\log^2 n}{B} + \frac{\log n}{\log B} + \frac{\log k}{\log B})$ | Indirekce - Zbavení se kvadrátu v logaritmu: Vytvoříme fragmenty, ve kterých budou uložené prvky, velikost fragmentu $\Theta(\log n)$. Do původní struktury si ukládám pouze reprezentanty. Najdu reprezentanta, tím najdu fragmenty (stojí $O(\frac{\log n}{\log B})$), do kterých prvek vkládám, pokud přeteče, tak řeším rozdělením fragmentu na dva a ten druhý vkládám. Tím operace globální operace $Insert_G(fragment)$ a ta nastane velikost fragmentu $:= \log n$ krát méně často, tedy amortizovaně proběhne při každém $Insertu$ $\frac{1}{\log n}$ operací globálního $Insertu_G$. Fragment je souvislá část paměti, operace s nimi v $O(\frac{\log n}{B})$. Amortizovaně jeden $Insert$ bude v $O(\frac{\log n}{B} + 1)$ I/O složitosti. $Find$ bude v $O(\frac{\log n}{\log B})$ I/O složitosti. | USPORNE DS: C je mn. objektů, min. počet bitů na uložení je OPT $=\lceil{\log|C|}\rceil$. Implicitní, Úsporné, Kompaktní DS. Reprezentace řetězce nad abecedou $\Sigma$. $\Sigma = [10] \rightarrow \lceil\log10\rceil = 4$ bity. Pokud budeme kódovat po dvojici, tak $\lceil\log_2{100}\rceil = 7$. Obecně pro $k$ znaků můžeme uložit do $k \cdot \log_2{10}+1$ bitů, lepší než $k \cdot (\log_2{10}+1)$. ALE zrušíme si lokální dekódovatelnost. $\text{Redundance} := \text{Skutečně zabraný prostor} - OPT$ | SOLE - Bezprefixové kódy  Zprávu na bloky velikosti $b$. Padding $10000\ldots$na konec zprávy, aby byla dělitelná $b$. Za každý blok bit $0$ nebo $1$, zda zpráva pokračuje, či končí. Redundance $= b + \lceil\frac{n}{b}\rceil = b + \frac{n}{b} + 1$.

\end{multicols*}
\end{document}